% 16 pages
% Structure: contributions (which absorb the representation question) ->
% the two remaining recurring considerations from \secref{sec:objectives} as
% sections -> limitations -> trend/outlook.
% Clinical and methodological impact are folded into the three threads rather
% than standing as separate sections (departs from Christian ch.7 deliberately).
% NOTE: if the faculty expects an explicit clinical-relevance section, the
% material for it is the availability + privacy + robustness bullets under
% "Where the geometry comes from" below.

% ---------------------------------------------------------------------------
% CHRIS'S ORIGINAL NOTES (preserved verbatim)
%
% On shape-generation (its getting much easier to generate mensingful shapes):
%   nnunet, Totalsegementator, cheap destilling -> \cite{grossbroehmer2026frozen}
% Maybe lets do some statistics on one random Totalsegmentator dataset:
%   How many voxels are labelled etc?
% On generalised medical image registration (its an active topic and other
%   appraoches): ICON, GradICON etc.
% ---------------------------------------------------------------------------

\section{Contributions}
% Chapter-by-chapter recap, with the representation question folded in here
% rather than given its own section (former thread 1).
%
% WHY FOLDED: a cross-chapter verdict on which representation is best cannot be
% supported. The five chapters use different representations on different tasks,
% anatomies and modalities, and were never controlled against one another. Say
% that limit plainly; it is stronger than implying a verdict.
%
% What CAN be said here, and what each claim rests on:
%  - The two genuine head-to-head comparisons in the thesis:
%      * ch:MIUA24, learned label encodings vs classical Euclidean Distance Maps
%      * ch:BVM25, INR-SSM vs linear PDM-SSM (partial rings; disc-to-ring, where
%        the topologies differ and one keypoint set cannot cover both)
%  - Analytical claims from the taxonomy (\secref{sec:bg-shapes}), which hold
%    without controlled experiments:
%      * an INR and a voxel-stored SDF share codomain and recovery rule and
%        differ only in how the field is stored
%      * representation 9 (deformation field) is not a family but the
%        denominator: it is what registration produces
%  - Which representation each chapter used and why that problem called for it:
%    learned encodings (ch:MIUA24), implicit shape model (ch:BVM25), continuous
%    field (ch:WBIR24), semantic label features (ch:TMI), point clouds
%    (ch:IJCARS).
%  - ch:WBIR24 sits differently from the rest: its link is representation
%    continuity (continuous vs discretised) rather than shape as such.
%
% SIGNPOSTING: \secref{sec:objectives} promises three recurring considerations
% and this chapter returns to two of them as sections. Say explicitly here that
% the first is answered in this recap, so the reader is not left waiting for a
% section that never arrives.
%
%  - Impact of ch:TMI beyond this thesis (uptake verified in the full text,
%    2026-10-08). Uptake:
%    Applications, where convexAdam was used to solve someone else's problem:
%      * CT-to-MRI registration that produced the aligned ground truth of the
%        public DoseRAD2026 challenge dataset (115 patients; hyperparameters
%        tuned by hand, not self-configured) \cite{xiao2026doserad}
%      * clinical HCC ablation-margin pipeline, 139 lesions: instance
%        optimisation around each lesion after a global network registration;
%        the resulting margin predicts local tumour progression
%        \cite{fogarollo2026ablative}
%      * T2/DWI alignment in a 24-hospital prostate MRI annotation pipeline,
%        +1.0 AUROC points on PROMIS \cite{bosma2025scale}. NOTE: Heinrich and
%        Hering are co-authors, so not fully independent uptake.
%    Methods that reuse the decoupled design and swap in their own features:
%      * DINO-Reg: DINOv2 features + convexAdam optimisation \cite{song2024dinoreg}
%      * benchmark of five vision foundation models (SAM, MedSAM, DINOv2, ...)
%        for breast MRI registration, all plugged into the convexAdam two-stage
%        optimisation; convexAdam is the fixed back-end that makes the
%        encoders comparable \cite{gu2025vfmready}
%      * anatomix (ICLR 2025): general-purpose features from randomised
%        synthesis, registered with "the ConvexAdam registration solver";
%        adding its features to convexAdam's handcrafted ones gains 11 and 6
%        Dice points (L2R AbdomenMRCT, MM-WHS) \cite{dey2025anatomix}
%      * HAGR (MICCAI 2026): learned modality-agnostic features, registered
%        with the convexAdam coupled convex scheme at each pyramid level plus
%        a final instance optimisation; convexAdam is also its strongest
%        non-learned baseline on Learn2Reg AbdomenMRCT \cite{mi2026hagr}
%      * SAMConvex: SAM embeddings + a cost-volume pyramid added to the
%        convexAdam convex step \cite{li2023samconvex}
%    For the thread section: the split into features and optimisation is
%    what lets others plug in new features or apply it to new tasks without
%    retraining anything.
%    Translation:
%      * convexAdam is a core part of pending patent applications in the EU
%        and the US; I am a co-inventor. Application numbers deliberately
%        left out of the public thesis.
%      * basis of the university spin-off CONVEA, currently being founded
%        under the EXIST programme; I am a co-founder.
%      * conflict of interest declared in the ch:TMI opening remarks.

\section{Where learning is used}
% Thread 2 from \secref{sec:objectives}. First of the two thread sections.
%  - The span: encodings trained to transfer across anatomical structures
%    without retraining (ch:MIUA24); a framework that keeps learning to a
%    minimum and configures itself per task instead (ch:TMI); representations
%    optimised per instance at test time rather than predicted by a trained
%    network (ch:BVM25, ch:WBIR24). (Old wording "no population model at all"
%    was wrong: both build population models, see the thesis map.)
%  - Connects to the universal-registration discussion in the outlook below:
%    minimising task-specific learning is what "universal" actually means in
%    practice.
%  - Discussion point: why the success of deep learning in segmentation did
%    not carry over to registration. Thesis to argue: DL wins the
%    representation (semantic features), optimisation keeps the deformation
%    (convex + Adam IO, INR fitting).
%      * no ground-truth deformations: unsupervised nets minimise the same
%        similarity + regulariser objective as classical optimisation, so
%        they amortise it (speed) rather than beat it (accuracy)
%      * amortised predictors break under domain shift, per-instance
%        optimisation adapts; hence IO bolted onto learned methods
%        (VoxelMorph++, uniGradICON+IO)
%      * the hard part is the similarity, i.e. the features: exactly where
%        DL helps (multimodal, semantic correspondence)
%      * evidence: independent work converges on DL features + convexAdam
%        solver (DINO-Reg, anatomix, HAGR, VFM benchmark; see Contributions)
%      * links the chapters: ch:TMI (learned features + optimisation),
%        ch:MIUA24 (learned shape encodings), ch:BVM25 / ch:WBIR24
%        (per-instance INR optimisation instead of a trained predictor);
%        and the Trend section (registration's shrinking share of the main venues)
%      * COUNTER-POSITION, discuss explicitly: uniGradICON
%        \cite{tian2024unigradicon}, a "foundation model" for registration.
%        One GradICON network trained on twelve public datasets (lung, knee,
%        brain, abdomen; CT, CBCT, MRI), claims in-distribution accuracy on
%        par with task-specific SOTA, zero-shot use on unseen tasks, and a
%        strong initialisation for finetuning. This is the opposite bet:
%        amortise the deformation itself across tasks instead of keeping
%        optimisation.
%        Points for the discussion (all from the papers, verified):
%          - its own limitations section names our recipe as the way
%            forward: multimodal support "could likely be improved by using
%            self-supervised modality-agnostic representations", and
%            "further improvements might also be possible by including
%            segmentations for training and instance optimization"
%          - it ships an IO variant (uniGradICON(IO)) itself
%          - zero-shot performance on unseen anatomical regions is "limited"
%            (their Tab. 4)
%          - anatomix \cite{dey2025anatomix} reports that uniGradICON+IO
%            "cannot handle the larger deformations in L2RAb", while their
%            features + convexAdam solver do
%        Its multimodal successor multiGradICON \cite{demir2024multigradicon}
%        (WBIR 2024) extends the bet to multimodal pairs (randomised
%        similarity loss, multimodal training data). Verified points:
%          - motivated by uniGradICON failing on large appearance gaps: it
%            "cannot register pairs containing T2w images, even with
%            instance optimization"
%          - its instance optimisation uses MIND-SSC, i.e. a handcrafted
%            descriptor from the classical toolbox (same family as the MIND
%            features convexAdam falls back on)
%          - remaining failure (DIXON fat/water) to be addressed "by
%            incorporating semantic information": again our recipe
%        Fair framing: uni-/multiGradICON win on speed and zero-setup use; the
%        open question is whether amortisation can match per-instance
%        optimisation on large deformations and new modalities.
%      * CAVEAT: do not claim DL registration failed. It won on speed and is
%        widely used; foundation registration models (uniGradICON, see above) are a
%        live counter-trend. Defensible claim: "DL won the representation,
%        optimisation kept the deformation".
%      * TODO: when the Learn2Reg 2026 solution papers are released, count
%        end-to-end DL vs DL features + optimisation/IO vs classical among the
%        entries and top ranks. Reframe if end-to-end dominates.
%      * TODO: check what \cite{hering2022learn2reg} reports on IO among the
%        top entries (not verified yet).
%  - Evidence from outside the thesis: the uptake of ch:TMI (others plugging
%    their own features into its solver) is listed under Contributions;
%    refer back to it here rather than repeating it.

\section{Where the geometry comes from}
% Thread 3 from \secref{sec:objectives}. The strongest of the three: it is
% argued in four of the five chapters rather than imposed. Absorbs the former
% "clinical impact" material.
%  - Availability is the enabling condition. Segmentation became cheap
%    (nnU-Net, TotalSegmentator, MONAI Model Zoo, and cheap distilling via a
%    frozen backbone plus a small head \cite{xu2025generalizable}), which
%    is *why* shape-as-representation became deployable at all. Argued outright
%    in ch:MIUA24; ch:TMI's discussion ends on the same trend.
%  - Chris's idea, worth doing: statistics on a TotalSegmentator dataset (how
%    many voxels labelled, how many structures, at what cost). Turns the
%    availability claim from an assertion into a number, and is far cheaper
%    than the publication-trend analysis.
%  - Dependence as a design axis, not an assumption: ch:TMI hedges across it,
%    using hand-crafted MIND where no annotations exist and nnU-Net semantic
%    features where they do, rather than committing to either.
%  - Degradation of the source: partial and non-overlapping sweeps
%    (ch:IJCARS), corrupted and incomplete segmentations (ch:BVM25).
%  - Correspondence as a special case: ch:BVM25 builds an SSM from
%    segmentation maps and SDFs without any direct correspondence, removing an
%    annotation requirement rather than automating it.
%  - Privacy: a method consuming only label maps never touches identifiable
%    intensity data, so it is privacy-preserving by design (ch:MIUA24). A real
%    clinical argument, currently buried in that chapter.

\section{Limitations}
%  - ch:WBIR24 does not participate in thread 3 at all: SINA derives its
%    representation from intensities and has no external geometry source.
%  - Semantic features still presuppose that annotations existed to train the
%    segmenter. Where no meaningful labels exist (CuRIOUS, NLST, LungCT in
%    ch:TMI) the method falls back to hand-crafted descriptors.
%    NOTE: \cite{grossbroehmer2026frozen} weakens this limitation considerably.
%  - Evaluation circularity: methods consuming segmentations must not be judged
%    on overlap alone (see \secref{sec:bg-registration-evaluation}).
%  - Scope: X-ray and ultrasound-to-CT work was excluded from this thesis.

\section{Trend}
% Chris's original note: "We might want to do a crawl of publications to assess
% the popularity growth? We have to do the crawl then." -> the crawl was done.
% Scope was cut to the level-1 finding only: registration is a declining share
% of the main venues, within venue as well as in aggregate. See
% scripts/litsearch/PLAN.md for the 13 Scopus queries and the numbers so far.

\section{Research Outlook}
% - New approaches in universal medical image registration. Chris's original
%   note: "or other approaches? image foundation models, domain adaptation, ?"
%   plus ICON, GradICON. Ties back to "Where learning is used".
% - Foundation segmentation models as the direction that makes geometry
%   universally available, and what that would change for all five chapters.
%   \cite{xu2025generalizable} shows that a frozen backbone plus a small head suffices.
% - Population models for registration plausibility. ch:BVM25 (robust SSM)
%   and ch:WBIR24 (sharp atlas) both BUILD population models, but no chapter
%   uses one to assess or constrain a registration. Natural next step: use
%   them to judge whether an estimated transformation is anatomically
%   plausible (e.g. shape of the warped structure under the SSM, or
%   deviation from atlas-derived deformation statistics), or as a prior
%   during optimisation. Closes the thread opened in the Motivation
%   ("combined with population-based information ... assess whether an
%   estimated transformation is anatomically plausible").
%   Also the link that places ch:WBIR24 in the shape argument.
%   Keep it short, two citations: population statistics of deformations
%   \cite{rueckert2003automatic} and a learned probabilistic deformation
%   model whose latent likelihood could serve as a plausibility score
%   \cite{krebs2019learning}. (Longer list from the lit search, not used:
%   Pennec 2005 Riemannian elasticity, Lüthi 2018 GPMM, Ehrhardt 2011 4D
%   lung motion model [Lübeck], Mansilla 2020 AC-RegNet, Sokooti 2019 error
%   prediction.)
% - TODO: related-work material to be MOVED here from background sec 2.3, where
%   it currently reads as an unfamiliar list; here it can be comparative against
%   the chapters:
%     "Shape as network input"         (Su & Yang 2023; Chen et al. 2021; Hu et al. 2019)
%     "Implicit shape representations" (Sun et al. 2022; Dalca et al. 2019)
%   Forward \chapref references in that text need reversing, and the
%   "% TODO primary cite" markers travel with it.
